% ============================================================================
%  অধ্যায় ৮ : বীজগণিত ২
%  COMPILE WITH XeLaTeX  (twice, so the point total resolves)
%  Overleaf: Menu -> Compiler -> XeLaTeX
%  Requires: bnhandout.sty  and  fonts/Kalpurush.ttf  in the project
% ============================================================================
\documentclass[11pt]{article}
\usepackage{bnhandout}

\usepackage{amsmath}
\usepackage{amssymb}
\usepackage{array}

\hauthor[https://jugantordey.github.io/]{Jugantor Dey}
\hseries{\textit{Mathematics for the Physics Olympiad}}

\begin{document}

\handouttitle{Chapter 8 : Algebra II}

\noindent
এই অধ্যায়ে ঢোকার আগে অধ্যায় ২ (বীজগণিত ১), ৩ (উৎপাদক বিশ্লেষণ) ও ৪
(ফাংশন ১) শেষ করে নাও। ২ থেকে লাগবে পূর্ণসংখ্যা সূচকের নিয়ম, ৩ থেকে লাগবে
মধ্যপদ বিশ্লেষণ ও দ্বিঘাত সমীকরণ, আর ৪ থেকে লাগবে function, domain, এবং সবচেয়ে
জরুরি, বিপরীত ফাংশন কখন থাকে সেই শর্তটি।

এটি Intermediate সেটের প্রথম অধ্যায়, আর এখান থেকে গতি একটু বাড়বে। তবে
বিষয়বস্তু নতুন হলেও পদ্ধতি পুরোনোই: প্রতিটি নিয়ম কোথা থেকে এল, সেটা আগে
দেখানো হবে, তারপর ব্যবহার।

একটা কথা সততার খাতিরে আগেই বলে রাখি। এই অধ্যায়ে তিনটি জায়গায় এমন দাবি
ব্যবহার করা হয়েছে যেগুলো এখানে প্রমাণ করা যায়নি, কারণ তাদের প্রমাণের জন্য
limit দরকার, আর limit আসবে অধ্যায় ১৪-তে। তিনটিই যথাস্থানে স্পষ্ট করে চিহ্নিত
করা আছে, আর কোনটির প্রমাণ কোথায় আসবে তাও বলা আছে। ধার করা জিনিস ধার করা বলেই
চেনানো ভালো; চুপচাপ চালিয়ে দিলে পরে বিশ্বাসটাই নড়বড়ে হয়ে যায়।

এই অধ্যায়ে মোট \totalpoints\ পয়েন্ট আছে।

% ============================================================================
\section{ভগ্নাংশ সূচক}
% ============================================================================

অধ্যায় ২-এ সূচকের নিয়মগুলো প্রমাণ করা হয়েছিল, তবে কেবল ধনাত্মক পূর্ণসংখ্যা
সূচকের জন্য। কারণ তখন $a^{n}$-এর অর্থই ছিল ``$a$-কে $n$ বার গুণ করা'', আর
$a^{1/2}$ বা $a^{-3}$ লেখাটার কোনো অর্থ ছিল না। আধখানা বার গুণ করা যায় না।

এই অনুচ্ছেদের কাজ ওই অর্থটাকে বড় করা। কিন্তু নতুন অর্থ যেমন খুশি বসানো যাবে না;
একটি নির্দিষ্ট নীতি মেনে বসাতে হবে।

\begin{idea}[সম্প্রসারণের নীতি]
নতুন সূচকের অর্থ এমনভাবে ঠিক করব, যাতে পুরোনো নিয়মগুলো অটুট থাকে। বিশেষ করে
\[
  a^{m} \cdot a^{n} = a^{m+n}, \qquad
  \frac{a^{m}}{a^{n}} = a^{m-n}, \qquad
  (a^{m})^{n} = a^{mn}
\]
নিয়ম তিনটি যেন সব সূচকের জন্যই সত্য থাকে।

এটি স্বেচ্ছাচার নয়, বরং উল্টোটা: নিয়মগুলোই নতুন অর্থটিকে বাধ্য করে দেয়, আর
নিচে দেখা যাবে প্রতিটি ক্ষেত্রে একটিমাত্র সম্ভাব্য অর্থই টিকে থাকে।
\end{idea}

\noindent
\textbf{শূন্য সূচক।} প্রথম নিয়মে $n = 0$ বসাই:
\[
  a^{m} \cdot a^{0} = a^{m+0} = a^{m} .
\]
$a \neq 0$ হলে দুই পাশকে $a^{m}$ দিয়ে ভাগ করা যায়, আর তখন $a^{0} = 1$ ছাড়া আর
কোনো সম্ভাবনাই থাকে না।

\textbf{ঋণাত্মক সূচক।} এবার $m = -n$ বসাই:
\[
  a^{n} \cdot a^{-n} = a^{0} = 1
  \qquad \Longrightarrow \qquad
  a^{-n} = \frac{1}{a^{n}} .
\]
আবারও একটিমাত্র উত্তর।

\textbf{একক ভগ্নাংশ সূচক।} তৃতীয় নিয়মে $m = \tfrac{1}{n}$ বসাই:
\[
  \left( a^{1/n} \right)^{n} = a^{\frac{1}{n} \cdot n} = a^{1} = a .
\]
অর্থাৎ $a^{1/n}$ এমন একটি সংখ্যা যার $n$ ঘাত $a$; সংক্ষেপে $a$-এর $n$-তম মূল।

\textbf{সাধারণ ভগ্নাংশ সূচক।} শেষে
\[
  a^{m/n} = \left( a^{1/n} \right)^{m} = \left( \sqrt[n]{a} \right)^{m} .
\]

\begin{idea}[ভিত্তি ধনাত্মক রাখতে হবে]
উপরের সবকিছুতে ধরে নেওয়া হয়েছে $a > 0$। এই শর্তটি বাদ দিলে নিয়মগুলোই ভেঙে
পড়ে, তাই এই সিরিজে ভগ্নাংশ সূচকের ভিত্তি সবসময় ধনাত্মক।
\end{idea}

\begin{remark}
শর্তটি কেন দরকার, তা একটি উদাহরণেই পরিষ্কার হয়। ধরা যাক $a = -8$ বসাতে দিলাম।
\[
  (-8)^{1/3} = \sqrt[3]{-8} = -2 ,
\]
যা নিরীহ মনে হয়। কিন্তু $\tfrac{1}{3} = \tfrac{2}{6}$, তাই একই সংখ্যাকে অন্যভাবে
লিখলে
\[
  (-8)^{2/6} = \left( (-8)^{2} \right)^{1/6} = 64^{1/6} = 2 .
\]
একই রাশির দুটি ভিন্ন মান, $-2$ আর $2$। অর্থাৎ ঋণাত্মক ভিত্তির ক্ষেত্রে
$a^{m/n}$ ভগ্নাংশটির লেখার উপর নির্ভর করে ফেলে, অথচ $\tfrac{1}{3}$ ও
$\tfrac{2}{6}$ একই সংখ্যা। অধ্যায় ৪-এর ভাষায় এটি ঠিক সেই ``দোটানা'', যা
function-এর সংজ্ঞা ভাঙে।

তাই $\sqrt[3]{-8} = -2$ লেখা চলবে, কারণ ওটা মূলের কথা; কিন্তু $(-8)^{1/3}$
লেখা চলবে না, কারণ ওটা সূচকের কথা, আর সূচকের নিয়ম সেখানে খাটে না।
\end{remark}

\begin{example}
সরল করো।
\[
  \text{(ক)}\ \ 32^{3/5}
  \qquad
  \text{(খ)}\ \ 81^{-3/4}
  \qquad
  \text{(গ)}\ \ \frac{x^{2/3} \cdot x^{5/6}}{x^{1/2}}
\]
\end{example}

\begin{solbox}
(ক) আগে মূল, পরে ঘাত নিলে সংখ্যাগুলো ছোট থাকে:
\[
  32^{3/5} = \left( 32^{1/5} \right)^{3} = 2^{3} = 8 .
\]

(খ) ঋণাত্মক সূচক মানে অন্যোন্যক:
\[
  81^{-3/4} = \frac{1}{81^{3/4}}
            = \frac{1}{\left( 81^{1/4} \right)^{3}}
            = \frac{1}{3^{3}} = \frac{1}{27} .
\]

(গ) সূচকগুলো নিয়ে কাজ করি, লবে যোগ আর ভাগে বিয়োগ:
\[
  \frac{2}{3} + \frac{5}{6} - \frac{1}{2}
  = \frac{4 + 5 - 3}{6} = \frac{6}{6} = 1 ,
\]
তাই রাশিটি $x^{1} = x$। এখানে $x > 0$ ধরে নেওয়া হয়েছে।
\end{solbox}

% ============================================================================
\section{লগারিদম}
% ============================================================================

\begin{idea}[লগারিদম মানে সূচকের বিপরীত]
$a > 0$ এবং $a \neq 1$ ধরো, আর বিবেচনা করো
\[
  f \colon \mathbb{R} \to (0, \infty), \qquad f(x) = a^{x} .
\]
এই $f$ একটি bijection। অধ্যায় ৪ অনুযায়ী তাই তার একটি inverse আছে, আর
সেটিকেই বলে $a$ ভিত্তিক logarithm:
\[
  \log_a \colon (0, \infty) \to \mathbb{R} ,
  \qquad
  \log_a x = y \iff a^{y} = x .
\]
সংজ্ঞাটি এক বাক্যে: $\log_a x$ হলো সেই সূচক, যা $a$-এর উপর বসালে $x$ পাওয়া যায়।

Inverse হওয়ার সরাসরি ফল দুটি বাতিলকরণ সমীকরণ:
\[
  a^{\log_a x} = x \quad (x > 0), \qquad \log_a (a^{y}) = y \quad (y \in \mathbb{R}) .
\]
\end{idea}

\begin{remark}
উপরে ``এই $f$ একটি bijection'' কথাটা দাবি হিসেবে রাখা হলো, প্রমাণ ছাড়া। এটি
এই অধ্যায়ের প্রথম ধার করা দাবি, তাই একটু খোলাসা করা যাক।

Injective হওয়ার অর্থ $a^{x}$ কখনো একই মান দুবার নেয় না, যা আসলে বলছে ফাংশনটি
সবসময় বাড়ে অথবা সবসময় কমে। Surjective হওয়ার অর্থ প্রতিটি ধনাত্মক সংখ্যাই
কোনো না কোনো $a^{x}$ আকারে লেখা যায়, অর্থাৎ ফাংশনটি $(0,\infty)$-এর কোনো মান
বাদ দেয় না।

দুটির কোনোটিই ভগ্নাংশ সূচকের সংজ্ঞা থেকে বেরোয় না, কারণ আমরা এখনো কেবল মূলদ
সূচকের অর্থ দিয়েছি; অমূলদ সূচকের অর্থই এখনো ঠিক করা হয়নি। এদের প্রমাণের জন্য
continuity লাগে, আর সেটি আসবে অধ্যায় ১৪-তে। ততক্ষণ আমরা দাবি দুটি মেনে
নিয়ে এগোচ্ছি, আর ছবির দিক থেকে এরা যথেষ্ট বিশ্বাসযোগ্যও।
\end{remark}

\begin{idea}[লগারিদমের নিয়ম]
$x, y > 0$ এবং $r$ যেকোনো বাস্তব সংখ্যা হলে
\[
  \log_a(xy) = \log_a x + \log_a y ,
\]
\[
  \log_a\!\left( \frac{x}{y} \right) = \log_a x - \log_a y ,
\]
\[
  \log_a (x^{r}) = r \log_a x .
\]
\end{idea}

\noindent
\textbf{প্রমাণ।} তিনটিরই একই কৌশল: লগারিদমগুলোকে নাম দিয়ে সূচকের ভাষায় নেমে
যাওয়া।

ধরি $u = \log_a x$ এবং $v = \log_a y$। সংজ্ঞা অনুযায়ী তখন $x = a^{u}$ এবং
$y = a^{v}$।

গুণ করলে, সূচকের প্রথম নিয়মে
\[
  xy = a^{u} \cdot a^{v} = a^{u+v} .
\]
সংজ্ঞাটি উল্টো দিকে পড়লে এর মানে $\log_a(xy) = u + v$, যা প্রথম দাবি।

ভাগ করলে $\dfrac{x}{y} = a^{u-v}$, যা দ্বিতীয় দাবি দেয়।

আর ঘাত নিলে $x^{r} = (a^{u})^{r} = a^{ur}$, যা তৃতীয় দাবি দেয়।
$\blacksquare$

\begin{idea}[ভিত্তি পরিবর্তন]
$a, b$ দুটিই ধনাত্মক ও $1$ নয় হলে
\[
  \log_b x = \frac{\log_a x}{\log_a b} .
\]
\end{idea}

\noindent
\textbf{প্রমাণ।} ধরি $y = \log_b x$, অর্থাৎ $b^{y} = x$। দুই পাশে $a$ ভিত্তিক
logarithm নিয়ে, তৃতীয় নিয়ম প্রয়োগ করে
\[
  y \log_a b = \log_a x .
\]
$b \neq 1$ বলে $\log_a b \neq 0$, তাই ভাগ করা বৈধ। $\blacksquare$

\begin{example}
সমাধান করো: $\ \log_3 x + \log_3 (x - 6) = 3$।
\end{example}

\begin{solbox}
প্রথম নিয়মে বাঁ পাশ জুড়ে দিই:
\[
  \log_3 \big( x(x-6) \big) = 3 .
\]
সংজ্ঞা অনুযায়ী এর মানে
\[
  x(x-6) = 3^{3} = 27
  \ \Longrightarrow\ x^{2} - 6x - 27 = 0 .
\]
অধ্যায় ৩-এর মধ্যপদ বিশ্লেষণে $(x-9)(x+3) = 0$, তাই $x = 9$ বা $x = -3$।

কিন্তু দুটোই গ্রহণ করা যাবে না। $\log_3$-এর domain হলো $(0,\infty)$, তাই মূল
সমীকরণে $x > 0$ এবং $x - 6 > 0$ দুটোই লাগে, অর্থাৎ $x > 6$। $x = -3$ ওই শর্ত
মানে না, ফলে সে বহিরাগত মূল।

একমাত্র সমাধান $x = 9$।
\end{solbox}

\begin{remark}
উপরের শেষ ধাপটা বাদ দিলে উত্তর ভুল হতো, আর ভুলটা ধরা পড়ত না, কারণ বীজগণিতে
কোথাও গোলমাল হয়নি। গোলমালটা domain-এ।

লক্ষ করো কী ঘটেছে। $\log x + \log(x-6)$ থেকে $\log(x(x-6))$-এ যাওয়ার সময়
domain চুপচাপ বড় হয়ে গেছে: বাঁ পাশের জন্য $x > 6$ লাগে, কিন্তু ডান পাশ
$x < 0$-তেও সংজ্ঞায়িত, কারণ তখন দুটি ঋণাত্মক সংখ্যার গুণফল ধনাত্মক। অধ্যায়
৪-এ ভগ্নাংশ কাটার সময় ঠিক এই একই ফাঁদ দেখা গিয়েছিল।

নিয়ম হিসেবে মনে রাখো: logarithm-এর সমীকরণ সমাধানের পর প্রতিটি উত্তর মূল
সমীকরণে বসিয়ে যাচাই করো।
\end{remark}

% ============================================================================
\section{সংখ্যা \texorpdfstring{$e$}{e}}
% ============================================================================

\begin{idea}[চক্রবৃদ্ধি সুদের প্রশ্ন]
ধরো ব্যাংকে $1$ টাকা রাখলে, আর বার্ষিক সুদের হার $100\%$। এক বছর শেষে পাবে
$2$ টাকা।

এবার ব্যাংক বলল, সুদ বছরে একবার নয়, দুবার দেওয়া হবে, প্রতিবার $50\%$ করে।
তাহলে ছয় মাস পরে $1.5$ টাকা, আর বছর শেষে $1.5 \times 1.5 = 2.25$ টাকা। আগের
চেয়ে বেশি।

বছরে $n$ বার দিলে প্রতিবার হার $\tfrac{1}{n}$, আর বছর শেষে
\[
  A_n = \left( 1 + \frac{1}{n} \right)^{n} .
\]
স্বাভাবিক প্রশ্ন: $n$ যত বড় করব, টাকা কি যত খুশি বাড়বে?
\end{idea}

\noindent
কয়েকটি মান হিসাব করে দেখা যাক।

\begin{center}
\begin{tabular}{@{}r l@{}}
$n$ & $\left( 1 + \tfrac{1}{n} \right)^{n}$ \\[2pt]
\hline \\[-8pt]
$1$ & $2$ \\
$2$ & $2.25$ \\
$4$ & $2.44140\ldots$ \\
$12$ & $2.61303\ldots$ \\
$365$ & $2.71456\ldots$ \\
$1000$ & $2.71692\ldots$ \\
$10000$ & $2.71814\ldots$ \\
\end{tabular}
\end{center}

\noindent
উত্তরটি ``না''। সংখ্যাগুলো বাড়ছে ঠিকই, কিন্তু ক্রমশ ধীরে, আর $2.72$ পেরোচ্ছে
না। মনে হচ্ছে তারা একটি নির্দিষ্ট সংখ্যার দিকে এগোচ্ছে।

\begin{idea}[সংখ্যা $e$]
$\left( 1 + \tfrac{1}{n} \right)^{n}$ রাশিটি $n$ বাড়ার সঙ্গে সঙ্গে একটি
নির্দিষ্ট সংখ্যার দিকে এগোয়। ওই সংখ্যাটিকে বলে $e$, আর
\[
  e = 2.718281828\ldots
\]
$e$ অমূলদ, তাই এর দশমিক বিস্তার কখনো শেষ হয় না ও পুনরাবৃত্তিও হয় না।

$e$ ভিত্তিক logarithm-কে বলে natural logarithm, আর তাকে $\ln$ লেখা হয়:
\[
  \ln x = \log_e x .
\]
\end{idea}

\begin{remark}
এখানে সৎ হওয়া দরকার, কারণ উপরের বাক্সে দুটি দাবি প্রমাণ ছাড়া বসানো হয়েছে।

প্রথমত, ``একটি নির্দিষ্ট সংখ্যার দিকে এগোয়'' কথাটা একটি limit-এর দাবি। উপরের
তালিকা দেখে সংখ্যাগুলো থিতিয়ে আসছে বলে মনে হচ্ছে বটে, কিন্তু মনে হওয়া আর প্রমাণ
এক জিনিস নয়। একটি ক্রম বাড়তে বাড়তেও অসীমে চলে যেতে পারে, কেবল ধীরে ধীরে।
তাই limit-টি আসলেই আছে কি না, সেটি প্রমাণ করতে হবে, আর সেই প্রমাণ আসবে
অধ্যায় ১৪-তে, যেখানে
\[
  e = \lim_{n \to \infty} \left( 1 + \frac{1}{n} \right)^{n}
\]
সমীকরণটি সংজ্ঞা হিসেবে লেখা হবে এবং তার বৈধতা দেখানো হবে।

দ্বিতীয়ত, $e$ অমূলদ হওয়ার দাবিটিও এখানে প্রমাণ করা হয়নি; ওটি একটি আলাদা
উপপাদ্য।

তাহলে এখানে $e$ নিয়ে কথা তোলা হলো কেন? কারণ পরের অধ্যায়গুলোতে তাকে বারবার
দরকার হবে, আর তার আসল পরিচয়টুকু আগে থেকে চেনা থাকলে সুবিধা। অধ্যায় ১০-এ
$e^{x}$ ও $\ln x$ পূর্ণ ফাংশন হিসেবে আসবে, ১৫-তে দেখা যাবে $e$-ই একমাত্র ভিত্তি
যার জন্য $\dfrac{d}{dx}a^{x} = a^{x}$, আর ১৮-এ অয়লারের সূত্রে সে জটিল সংখ্যার
সঙ্গে জুড়ে যাবে। আজ কেবল সংখ্যাটার সঙ্গে পরিচয়।
\end{remark}

\begin{example}
$10000$ টাকা বার্ষিক $8\%$ হারে $5$ বছরের জন্য রাখা হলো, আর সুদ ত্রৈমাসিক
ভিত্তিতে যোগ হয়। শেষে কত টাকা হবে? $e^{0.4} \approx 1.49182$ ধরে অবিরত
চক্রবৃদ্ধির ক্ষেত্রে উত্তরটি কী হতো, তুলনা করো।
\end{example}

\begin{solbox}
ত্রৈমাসিক মানে বছরে $4$ বার, তাই $5$ বছরে মোট $20$ বার সুদ বসবে, আর প্রতিবারের
হার $\tfrac{0.08}{4} = 0.02$। সুতরাং
\[
  A = 10000\,(1.02)^{20} .
\]
$(1.02)^{20} \approx 1.48595$, তাই $A \approx 14859.47$ টাকা।

অবিরত চক্রবৃদ্ধিতে সুদ প্রতি মুহূর্তে বসে, আর সূত্রটি দাঁড়ায়
$A = P e^{rt}$; এই সূত্রটি এখানে ধার করা, এর উৎপত্তি আসবে অধ্যায় ১৯-এ।
মান বসিয়ে
\[
  A = 10000\, e^{0.08 \times 5} = 10000\, e^{0.4} \approx 14918.25 \text{ টাকা} .
\]

পার্থক্য প্রায় $59$ টাকা, অর্থাৎ শতকরা আধ ভাগেরও কম। এটাই $e$-এর গল্পের
নৈতিক শিক্ষা: যত ঘন ঘনই সুদ বসাও, লাভ অসীম হয় না, একটি ছাদে গিয়ে থামে।
\end{solbox}

% ============================================================================
\section{অনুক্রম ও ধারা}
% ============================================================================

\begin{idea}[অনুক্রম ও ধারা]
অনুক্রম (sequence) হলো একটি নির্দিষ্ট ক্রমে সাজানো সংখ্যার তালিকা:
$a_1, a_2, a_3, \dots$

অধ্যায় ৪-এর ভাষায় এটি আসলে একটি function, যার domain $\mathbb{N}$:
\[
  a \colon \mathbb{N} \to \mathbb{R} , \qquad a(n) = a_n .
\]
লেখার ধরনটাই কেবল আলাদা; $a(n)$-এর বদলে $a_n$ লেখা হয়। নতুন কোনো বস্তু এখানে
তৈরি হয়নি।

ধারা (series) হলো অনুক্রমের পদগুলোর যোগফল: $a_1 + a_2 + \cdots + a_n$।
প্রথম $n$টি পদের যোগফলকে লেখা হয় $S_n$।
\end{idea}

\begin{idea}[সমান্তর ধারা]
যে অনুক্রমে পরপর দুই পদের অন্তর ধ্রুবক, তাকে বলে সমান্তর অনুক্রম, আর ওই ধ্রুবক
অন্তরটিকে বলে সাধারণ অন্তর $d$। প্রথম পদ $a$ হলে
\[
  a_n = a + (n-1)d ,
\]
আর প্রথম $n$টি পদের যোগফল
\[
  S_n = \frac{n}{2}\Big( 2a + (n-1)d \Big) = \frac{n}{2}(a + l) ,
\]
যেখানে $l = a_n$ হলো শেষ পদ।
\end{idea}

\noindent
\textbf{উৎপত্তি।} $a_n$-এর সূত্রটি সরাসরি: প্রথম পদ থেকে $n$-তম পদে যেতে $d$
যোগ করতে হয় $n-1$ বার।

যোগফলের সূত্রটির কৌশল সুন্দর। যোগফলটি একবার সোজা করে, আরেকবার উল্টো করে লিখি:
\[
  S_n = a + (a+d) + (a+2d) + \cdots + l ,
\]
\[
  S_n = l + (l-d) + (l-2d) + \cdots + a .
\]
এবার দুটি সমীকরণ পদে পদে যোগ করি। প্রথম স্তম্ভে পাই $a + l$; দ্বিতীয় স্তম্ভে
$(a+d) + (l-d) = a + l$; তৃতীয় স্তম্ভে $(a+2d) + (l-2d) = a + l$। প্রতিটি
স্তম্ভেই $d$-গুলো কেটে গিয়ে $a + l$ থাকছে, আর স্তম্ভ আছে $n$টি। সুতরাং
\[
  2 S_n = n(a + l)
  \qquad \Longrightarrow \qquad
  S_n = \frac{n}{2}(a+l) .
\]
এতে $l = a + (n-1)d$ বসালেই প্রথম রূপটি পাওয়া যায়। $\blacksquare$

\begin{idea}[গুণোত্তর ধারা]
যে অনুক্রমে পরপর দুই পদের অনুপাত ধ্রুবক, তাকে বলে গুণোত্তর অনুক্রম, আর ওই
ধ্রুবক অনুপাতটিকে বলে সাধারণ অনুপাত $r$। প্রথম পদ $a$ হলে
\[
  a_n = a r^{\,n-1} ,
\]
আর $r \neq 1$ হলে
\[
  S_n = \frac{a(1 - r^{n})}{1 - r} .
\]
\end{idea}

\noindent
\textbf{উৎপত্তি।} যোগফলটি লিখি, আর তাকে $r$ দিয়ে গুণ করে আরেকটি সমীকরণ বানাই:
\[
  S_n = a + ar + ar^{2} + \cdots + ar^{\,n-1} ,
\]
\[
  r S_n = \phantom{a + {}} ar + ar^{2} + \cdots + ar^{\,n-1} + ar^{n} .
\]
বিয়োগ করলে মাঝের সব পদ জোড়ায় জোড়ায় কেটে যায়, কেবল দুই প্রান্তের দুটি থাকে:
\[
  S_n - r S_n = a - ar^{n}
  \qquad \Longrightarrow \qquad
  S_n (1 - r) = a(1 - r^{n}) .
\]
$r \neq 1$ বলে ভাগ করা বৈধ। $r = 1$ হলে সব পদই $a$, তাই তখন সহজভাবে
$S_n = na$। $\blacksquare$

\begin{idea}[অসীম গুণোত্তর ধারা]
$|r| < 1$ হলে গুণোত্তর ধারাটির অসীম পর্যন্ত যোগফল অর্থপূর্ণ, আর
\[
  S = \frac{a}{1 - r} .
\]
\end{idea}

\noindent
\textbf{উৎপত্তি, একটি ধার করা ধাপ সহ।} উপরের সূত্রটি নিয়ে $n$ বড় করতে থাকি।
$|r| < 1$ হলে $r^{n}$ ক্রমশ ছোট হতে হতে $0$-এর দিকে এগোয়; যেমন
$r = \tfrac{1}{2}$ হলে $r^{10} \approx 0.001$ আর $r^{20} \approx 0.000001$।

\emph{এই ধাপটিই ধার করা।} ``$r^{n}$ শূন্যের দিকে এগোয়'' একটি limit-এর দাবি,
আর তার প্রমাণ আসবে অধ্যায় ১৪-তে। দাবিটি মেনে নিলে সূত্রে $r^{n}$-এর
জায়গায় $0$ বসিয়ে
\[
  S = \frac{a(1-0)}{1-r} = \frac{a}{1-r} .
\]

$|r| \geq 1$ হলে $r^{n}$ ছোট হয় না, তাই যোগফলও কোনো সংখ্যায় থিতোয় না। শর্তটি
তাই আলংকারিক নয়। $\blacksquare$

\begin{example}
একটি সমান্তর অনুক্রমের তৃতীয় পদ $11$ ও সপ্তম পদ $27$। প্রথম পদ, সাধারণ অন্তর
এবং প্রথম $20$টি পদের যোগফল বের করো।
\end{example}

\begin{solbox}
সূত্রে বসিয়ে দুটি সমীকরণ পাই:
\[
  a + 2d = 11 , \qquad a + 6d = 27 .
\]
অধ্যায় ২-এর অপনয়ন পদ্ধতিতে বিয়োগ করলে $4d = 16$, তাই $d = 4$, আর প্রথমটি
থেকে $a = 11 - 8 = 3$।

যোগফল:
\[
  S_{20} = \frac{20}{2}\Big( 2 \cdot 3 + 19 \cdot 4 \Big)
         = 10\,(6 + 76) = 820 .
\]
\end{solbox}

\begin{remark}
অসীম ধারার একটি বিখ্যাত ফল এখান থেকেই বেরোয়। ধরো
\[
  0.999\ldots = 0.9 + 0.09 + 0.009 + \cdots ,
\]
যা $a = 0.9$ ও $r = 0.1$ বিশিষ্ট গুণোত্তর ধারা। সূত্রে
\[
  S = \frac{0.9}{1 - 0.1} = \frac{0.9}{0.9} = 1 .
\]
অর্থাৎ $0.999\ldots$ আর $1$ একই সংখ্যা, ``প্রায় সমান'' নয়, হুবহু সমান।

অনেকেরই এটি মেনে নিতে আপত্তি হয়, আর আপত্তির উৎসটা বোঝা দরকার। $0.999\ldots$
লেখাটা কোনো প্রক্রিয়া নয় যা $1$-এর দিকে ছুটছে; ওটা একটি সংখ্যা, আর ওই সংখ্যাটির
সংজ্ঞাই হলো উপরের অসীম যোগফল। যোগফলটি $1$, তাই সংখ্যাটিও $1$।
\end{remark}

% ============================================================================
\section{সিগমা নোটেশন}
% ============================================================================

\begin{idea}[যোগচিহ্ন]
লম্বা যোগফল বারবার লেখা ক্লান্তিকর, তাই তার একটি সংক্ষিপ্ত লিপি আছে:
\[
  \sum_{i=1}^{n} a_i \;=\; a_1 + a_2 + \cdots + a_n .
\]
$\sum$ হলো গ্রিক বড় হাতের sigma। নিচে লেখা $i = 1$ বলে গণনা কোথা থেকে শুরু,
উপরে লেখা $n$ বলে কোথায় শেষ, আর $a_i$ বলে প্রতিটি পদ দেখতে কেমন।

$i$ চলকটির নাম নিছকই একটি নাম; $\sum_{i=1}^{n} a_i$ আর $\sum_{j=1}^{n} a_j$
একই সংখ্যা। একে বলে dummy variable।

গুণফলের জন্যও অনুরূপ চিহ্ন আছে, বড় হাতের pi:
\[
  \prod_{i=1}^{n} a_i \;=\; a_1 \cdot a_2 \cdots a_n .
\]
\end{idea}

\begin{remark}
এই চিহ্নটি দেখে ঘাবড়ানোর কিছু নেই। $\sum$ কোনো নতুন গাণিতিক ক্রিয়া নয়, নতুন
কোনো ক্ষমতাও দেয় না; সে নিছক হিসাবরক্ষণের একটি সংক্ষিপ্ত লিপি। যা লেখা হচ্ছে
তা সেই চেনা যোগই।

তবু লিপিটা শেখা জরুরি, কারণ সামনের অধ্যায়গুলো এতেই লেখা হবে। বিশেষ করে
অধ্যায় ১৭-এ integral-এর সংজ্ঞাই দেওয়া হবে একটি $\sum$ দিয়ে, যেখানে পদের
সংখ্যা অসীমের দিকে যায়।
\end{remark}

\begin{idea}[দুটি নিয়ম, আর একটি ফাঁদ]
যোগের ক্রম ও গুচ্ছবিন্যাস বদলানো যায়, তাই
\[
  \sum_{i=1}^{n} (c\,a_i) = c \sum_{i=1}^{n} a_i ,
  \qquad
  \sum_{i=1}^{n} (a_i + b_i) = \sum_{i=1}^{n} a_i + \sum_{i=1}^{n} b_i .
\]
প্রথমটি নিছক বণ্টন বিধি, দ্বিতীয়টি নিছক পদ পুনর্বিন্যাস। কোনোটিই গভীর কিছু নয়।

কিন্তু গুণের বেলায় অনুরূপ নিয়ম \emph{নেই}:
\[
  \sum_{i=1}^{n} (a_i b_i) \neq \left( \sum_{i=1}^{n} a_i \right)
                                 \left( \sum_{i=1}^{n} b_i \right) .
\]
$n = 2$ ধরে $a = (1,2)$, $b = (3,4)$ নিলে বাঁ পাশ $3 + 8 = 11$, ডান পাশ
$3 \times 7 = 21$। একটিমাত্র প্রতি-উদাহরণই যথেষ্ট।
\end{idea}

\begin{idea}[দুটি দরকারি যোগফল]
\[
  \sum_{i=1}^{n} i = \frac{n(n+1)}{2} ,
  \qquad
  \sum_{i=1}^{n} i^{2} = \frac{n(n+1)(2n+1)}{6} .
\]
\end{idea}

\noindent
\textbf{উৎপত্তি।} প্রথমটি একটি সমান্তর ধারা, যার $a = 1$, $l = n$, পদসংখ্যা
$n$। আগের অনুচ্ছেদের সূত্রে সরাসরি $S = \tfrac{n}{2}(1 + n)$।

দ্বিতীয়টির জন্য একটি কৌশল, যার নাম telescoping। লক্ষ করো
\[
  (k+1)^{3} - k^{3} = 3k^{2} + 3k + 1 .
\]
এবার $k = 1$ থেকে $n$ পর্যন্ত দুই পাশ যোগ করি। বাঁ পাশে পরপর পদগুলো কেটে যায়:
$(2^3 - 1^3) + (3^3 - 2^3) + \cdots + ((n+1)^3 - n^3)$, যেখানে $2^3$ থেকে $n^3$
পর্যন্ত প্রতিটি একবার যোগ ও একবার বিয়োগ হয়ে বিদায় নেয়। থাকে কেবল
$(n+1)^{3} - 1$।

ডান পাশে, উপরের দুটি নিয়ম ব্যবহার করে,
\[
  3 \sum_{i=1}^{n} i^{2} + 3 \sum_{i=1}^{n} i + n .
\]
$S_2 = \sum i^{2}$ লিখে দুই পাশ সমান বসাই:
\[
  (n+1)^{3} - 1 = 3 S_2 + 3 \cdot \frac{n(n+1)}{2} + n .
\]
বাঁ পাশ খুলে $n^{3} + 3n^{2} + 3n$। সব কিছু দুই দিয়ে গুণ করে গুছিয়ে লিখলে
\[
  6 S_2 = 2n^{3} + 6n^{2} + 6n - 3n(n+1) - 2n
        = 2n^{3} + 3n^{2} + n
        = n(n+1)(2n+1) ,
\]
সুতরাং $S_2 = \dfrac{n(n+1)(2n+1)}{6}$. $\blacksquare$

\begin{example}
$\displaystyle \sum_{i=1}^{n} (3i^{2} - 2i + 1)$ নির্ণয় করো, এবং $n = 10$-এর
জন্য মান বের করো।
\end{example}

\begin{solbox}
নিয়ম দুটি খাটিয়ে যোগফলটিকে তিন টুকরো করি:
\[
  \sum_{i=1}^{n} (3i^{2} - 2i + 1)
  = 3 \sum_{i=1}^{n} i^{2} \;-\; 2 \sum_{i=1}^{n} i \;+\; \sum_{i=1}^{n} 1 .
\]
শেষ যোগফলটিতে প্রতিটি পদ $1$, আর পদ আছে $n$টি, তাই সে $n$। বাকি দুটিতে সূত্র
বসাই:
\[
  = 3 \cdot \frac{n(n+1)(2n+1)}{6} - 2 \cdot \frac{n(n+1)}{2} + n
  = \frac{n(n+1)(2n+1)}{2} - n(n+1) + n .
\]
$n$ সাধারণ উৎপাদক নিয়ে গুছিয়ে লিখি:
\[
  = n \left[ \frac{(n+1)(2n+1)}{2} - (n+1) + 1 \right]
  = n \left[ \frac{2n^{2} + 3n + 1 - 2n}{2} \right]
  = \frac{n(2n^{2} + n + 1)}{2} .
\]

$n = 10$ বসিয়ে
\[
  \frac{10\,(200 + 10 + 1)}{2} = \frac{10 \times 211}{2} = 1055 .
\]
\end{solbox}

% ============================================================================
\section{দ্বিপদী উপপাদ্য}
% ============================================================================

$(a+b)^{n}$ রাশিটি খুলে লেখা ছোট $n$-এর জন্য হাতেই করা যায়:
\[
  (a+b)^{0} = 1 ,
\]
\[
  (a+b)^{1} = a + b ,
\]
\[
  (a+b)^{2} = a^{2} + 2ab + b^{2} ,
\]
\[
  (a+b)^{3} = a^{3} + 3a^{2}b + 3ab^{2} + b^{3} ,
\]
\[
  (a+b)^{4} = a^{4} + 4a^{3}b + 6a^{2}b^{2} + 4ab^{3} + b^{4} .
\]
$n = 10$-এর জন্য এভাবে করতে গেলে সারাদিন লেগে যাবে। তাই একটি নিয়ম খোঁজা যাক।

দুটি জিনিস চোখে পড়ে। প্রথমত, প্রতিটি পদে $a$ ও $b$-এর সূচকের যোগফল $n$, আর
$a$-এর সূচক $n$ থেকে $0$-তে নামে যখন $b$-এর সূচক $0$ থেকে $n$-এ ওঠে। দ্বিতীয়ত,
সহগগুলো একটি চেনা বিন্যাস তৈরি করে।

\begin{idea}[প্যাসকেলের ত্রিভুজ]
কেবল সহগগুলো সাজিয়ে লিখলে পাওয়া যায়
\[
\begin{array}{ccccccccccc}
  & & & & & 1 & & & & & \\[2pt]
  & & & & 1 & & 1 & & & & \\[2pt]
  & & & 1 & & 2 & & 1 & & & \\[2pt]
  & & 1 & & 3 & & 3 & & 1 & & \\[2pt]
  & 1 & & 4 & & 6 & & 4 & & 1 & \\[2pt]
  1 & & 5 & & 10 & & 10 & & 5 & & 1
\end{array}
\]
প্রতিটি সারির দুই প্রান্তে $1$, আর ভেতরের প্রতিটি সংখ্যা তার ঠিক উপরের দুটি
সংখ্যার যোগফল। যেমন শেষ সারির $10$ এসেছে তার উপরের $4$ ও $6$ থেকে।

এই নিয়মে যত ইচ্ছা সারি বাড়ানো যায়, আর $n$-তম সারিটিই $(a+b)^{n}$-এর সহগ।
\end{idea}

\begin{idea}[দ্বিপদী সহগ]
প্যাসকেলের ত্রিভুজ ছোট $n$-এর জন্য চমৎকার, কিন্তু $(a+b)^{40}$-এর একটিমাত্র
পদ চাইলে চল্লিশটি সারি লেখা অবাস্তব। তাই সহগের একটি সরাসরি সূত্র দরকার।

প্রথমে factorial: ধনাত্মক পূর্ণসংখ্যা $n$-এর জন্য
$n! = 1 \cdot 2 \cdot 3 \cdots n$, আর প্রথা অনুযায়ী $0! = 1$।

এবার $0 \leq k \leq n$ হলে দ্বিপদী সহগ
\[
  \binom{n}{k} = \frac{n!}{k!\,(n-k)!} .
\]
এটিই প্যাসকেলের ত্রিভুজের $n$-তম সারির $k$-তম সংখ্যা, যেখানে দুই দিকেই গণনা $0$
থেকে শুরু। যেমন
\[
  \binom{4}{2} = \frac{4!}{2!\,2!} = \frac{24}{4} = 6 ,
\]
যা সত্যিই $(a+b)^{4}$-এর মাঝের সহগ।
\end{idea}

\begin{idea}[প্যাসকেলের নিয়ম]
$1 \leq k \leq n-1$ হলে
\[
  \binom{n-1}{k-1} + \binom{n-1}{k} = \binom{n}{k} .
\]
এটিই ``উপরের দুটি যোগ করো'' নিয়মটির বীজগাণিতিক রূপ।
\end{idea}

\noindent
\textbf{প্রমাণ।} বাঁ পাশ লিখি সংজ্ঞা দিয়ে:
\[
  \frac{(n-1)!}{(k-1)!\,(n-k)!} + \frac{(n-1)!}{k!\,(n-1-k)!} .
\]
হর দুটিকে এক করতে প্রথম ভগ্নাংশের লব ও হরকে $k$ দিয়ে এবং দ্বিতীয়টির লব ও হরকে
$(n-k)$ দিয়ে গুণ করি। তখন $(k-1)! \cdot k = k!$ এবং
$(n-1-k)! \cdot (n-k) = (n-k)!$, তাই দুটিরই হর দাঁড়ায় $k!\,(n-k)!$:
\[
  = \frac{(n-1)!\,k}{k!\,(n-k)!} + \frac{(n-1)!\,(n-k)}{k!\,(n-k)!}
  = \frac{(n-1)!\,\big( k + n - k \big)}{k!\,(n-k)!}
  = \frac{(n-1)!\,n}{k!\,(n-k)!} .
\]
আর $(n-1)! \cdot n = n!$, তাই এটি $\dbinom{n}{k}$. $\blacksquare$

\begin{idea}[দ্বিপদী উপপাদ্য]
যেকোনো ধনাত্মক পূর্ণসংখ্যা $n$-এর জন্য
\[
  (a+b)^{n} = \sum_{k=0}^{n} \binom{n}{k} a^{\,n-k} b^{\,k} .
\]
\end{idea}

\begin{remark}
এই দাবিটি এখানে প্রমাণ করা হলো না, আর সেটা স্পষ্ট করে বলা দরকার।

উপরে প্যাসকেলের নিয়মটি প্রমাণ করা হয়েছে, অর্থাৎ দেখানো হয়েছে সহগগুলো ঠিক
ওইভাবেই একটি থেকে আরেকটি তৈরি হয়। কিন্তু ``প্রতিটি সারি সত্যিই $(a+b)^{n}$-এর
সহগ দেয়'' কথাটা এর চেয়ে বেশি; সেটি প্রমাণ করতে হলে $n$ থেকে $n+1$-এ যাওয়ার
ধাপটি বাঁধতে হয়, আর তার জন্য দরকার induction বা $\binom{n}{k}$-এর গণনামূলক
অর্থ। এই সিরিজে এখনো কোনোটিই আসেনি।

তাই আপাতত উপপাদ্যটি ব্যবহারের নিয়ম হিসেবে নেওয়া হচ্ছে, প্রমাণিত সত্য হিসেবে
নয়। ছোট $n$-এর জন্য তুমি নিজেই হাতে খুলে মিলিয়ে দেখতে পারো, আর নিচের উদাহরণে
$n = 4$-এর জন্য সেটাই করা হয়েছে।

একটি ভবিষ্যৎ ইঙ্গিতও দিয়ে রাখি। $n$ ধনাত্মক পূর্ণসংখ্যা না হলে, যেমন
$n = \tfrac{1}{2}$ বা $n = -1$, তখনও $(1+x)^{n}$-এর একটি বিস্তার আছে, তবে সেটি
আর শেষ হয় না; সে একটি অসীম ধারা। ওটি আসবে অধ্যায় ১৫-তে, Taylor সিরিজের
অংশ হিসেবে, আর পদার্থবিজ্ঞানে আসন্নীকরণের জন্য সেটিই সবচেয়ে বেশি ব্যবহৃত হয়।
\end{remark}

\begin{example}
$(2x - 3)^{4}$ বিস্তার করো।
\end{example}

\begin{solbox}
সূত্রে $a = 2x$, $b = -3$, $n = 4$ ধরি। সহগগুলো প্যাসকেলের চতুর্থ সারি:
$1, 4, 6, 4, 1$। পদগুলো একে একে:
\[
  k=0: \quad 1 \cdot (2x)^{4} = 16x^{4} ,
\]
\[
  k=1: \quad 4 \cdot (2x)^{3}(-3) = 4 \cdot 8x^{3} \cdot (-3) = -96x^{3} ,
\]
\[
  k=2: \quad 6 \cdot (2x)^{2}(-3)^{2} = 6 \cdot 4x^{2} \cdot 9 = 216x^{2} ,
\]
\[
  k=3: \quad 4 \cdot (2x)(-3)^{3} = 4 \cdot 2x \cdot (-27) = -216x ,
\]
\[
  k=4: \quad 1 \cdot (-3)^{4} = 81 .
\]
সুতরাং
\[
  (2x-3)^{4} = 16x^{4} - 96x^{3} + 216x^{2} - 216x + 81 .
\]
চিহ্নগুলো একবার পরপর বদলাচ্ছে, কারণ $b$ ঋণাত্মক আর তার সূচক এক এক করে বাড়ছে।

যাচাই করতে চাইলে $x = 1$ বসাও: বাঁ পাশ $(-1)^{4} = 1$, আর ডান পাশ
$16 - 96 + 216 - 216 + 81 = 1$। মিলে গেছে।
\end{solbox}

% ============================================================================
\section{সমস্যা}
% ============================================================================

\prob{1} সরল করো।
\begin{enumerate}
  \item $16^{3/4}$
  \item $27^{-2/3}$
  \item $\left( x^{1/2} y^{-1/3} \right)^{6}$, যেখানে $x, y > 0$
\end{enumerate}

\begin{sol}
\begin{enumerate}
  \item $\left( 16^{1/4} \right)^{3} = 2^{3} = 8$।
  \item $\dfrac{1}{27^{2/3}} = \dfrac{1}{\left( 27^{1/3} \right)^{2}}
        = \dfrac{1}{9}$।
  \item সূচকগুলো $6$ দিয়ে গুণ হয়:
        $x^{3} y^{-2} = \dfrac{x^{3}}{y^{2}}$।
\end{enumerate}
\end{sol}

\prob{2} সমাধান করো: $\ 9^{\,x+1} = 27^{\,x-1}$।

\begin{sol}
দুই পাশকে একই ভিত্তিতে আনি। $9 = 3^{2}$ ও $27 = 3^{3}$, তাই
\[
  3^{2(x+1)} = 3^{3(x-1)} .
\]
ভিত্তি এক এবং $3^{x}$ injective, তাই সূচক দুটি সমান:
\[
  2x + 2 = 3x - 3 \ \Longrightarrow\ x = 5 .
\]
যাচাই: বাঁ পাশ $9^{6} = 3^{12}$, ডান পাশ $27^{4} = 3^{12}$। মিলেছে।
\end{sol}

\prob{1} মান নির্ণয় করো।
\begin{enumerate}
  \item $\log_2 32$
  \item $\log_9 3$
  \item $\log_5 \dfrac{1}{25}$
\end{enumerate}

\begin{sol}
\begin{enumerate}
  \item $2^{5} = 32$, তাই $5$।
  \item $9^{1/2} = 3$, তাই $\tfrac{1}{2}$।
  \item $5^{-2} = \tfrac{1}{25}$, তাই $-2$।
\end{enumerate}
\end{sol}

\prob{2} $\log_{10} 2 = 0.3010$ ও $\log_{10} 3 = 0.4771$ ধরে নিচেরগুলো বের করো।
\begin{enumerate}
  \item $\log_{10} 12$
  \item $\log_{10} 1.5$
  \item $\log_{10} \sqrt{18}$
\end{enumerate}

\begin{sol}
\begin{enumerate}
  \item $12 = 2^{2} \cdot 3$, তাই
        $2(0.3010) + 0.4771 = 0.6020 + 0.4771 = 1.0791$।
  \item $1.5 = \tfrac{3}{2}$, তাই $0.4771 - 0.3010 = 0.1761$।
  \item $18 = 2 \cdot 3^{2}$, তাই $\log_{10} 18 = 0.3010 + 2(0.4771) = 1.2552$।
        বর্গমূল মানে সূচক $\tfrac{1}{2}$, তাই উত্তর
        $\tfrac{1}{2}(1.2552) = 0.6276$।
\end{enumerate}
\end{sol}

\prob{3} সমাধান করো: $\ \log_2 x + \log_2 (x - 2) = 3$।

\begin{sol}
বাঁ পাশ জুড়ে $\log_2 \big( x(x-2) \big) = 3$, অর্থাৎ
\[
  x^{2} - 2x = 8 \ \Longrightarrow\ x^{2} - 2x - 8 = 0
  \ \Longrightarrow\ (x-4)(x+2) = 0 ,
\]
তাই $x = 4$ বা $x = -2$।

Domain যাচাই করি। মূল সমীকরণে $\log_2 x$ থাকতে হলে $x > 0$, আর
$\log_2(x-2)$ থাকতে হলে $x > 2$; দুটো মিলে $x > 2$। $x = -2$ এই শর্ত মানে না,
তাই সে বহিরাগত মূল।

একমাত্র সমাধান $x = 4$। যাচাই: $\log_2 4 + \log_2 2 = 2 + 1 = 3$।
\end{sol}

\prob{2} $5^{x} = 12$ সমীকরণটি সমাধান করো। $\log_{10} 2 = 0.3010$ ও
$\log_{10} 3 = 0.4771$ ধরে নাও।

\begin{sol}
দুই পাশে $10$ ভিত্তিক logarithm নিয়ে, তৃতীয় নিয়মে
\[
  x \log_{10} 5 = \log_{10} 12 .
\]
$5 = \tfrac{10}{2}$, তাই $\log_{10} 5 = 1 - 0.3010 = 0.6990$। আর আগের সমস্যা
থেকে $\log_{10} 12 = 1.0791$। সুতরাং
\[
  x = \frac{1.0791}{0.6990} \approx 1.544 .
\]
যুক্তিসঙ্গত কি না দেখে নাও: $5^{1} = 5$ আর $5^{2} = 25$, তাই উত্তর $1$ ও $2$-এর
মাঝে থাকারই কথা।
\end{sol}

\prob{2} $5000$ টাকা বার্ষিক $6\%$ হারে $3$ বছরের জন্য রাখা হলো, সুদ
ষাণ্মাসিক ভিত্তিতে যোগ হয়। শেষে কত টাকা হবে? $(1.03)^{6} \approx 1.19405$
ধরে নাও।

\begin{sol}
ষাণ্মাসিক মানে বছরে $2$ বার, তাই $3$ বছরে $6$ বার সুদ বসবে, আর প্রতিবারের হার
$\tfrac{0.06}{2} = 0.03$। সুতরাং
\[
  A = 5000\,(1.03)^{6} \approx 5000 \times 1.19405 = 5970.25 \text{ টাকা} .
\]
সরল সুদে হতো $5000 + 5000(0.06)(3) = 5900$ টাকা, অর্থাৎ চক্রবৃদ্ধিতে প্রায়
$70$ টাকা বেশি।
\end{sol}

\prob{2} একটি সমান্তর অনুক্রমের প্রথম পদ $7$ ও সাধারণ অন্তর $5$।
\begin{enumerate}
  \item $25$-তম পদ বের করো।
  \item প্রথম $25$টি পদের যোগফল বের করো।
  \item কোন পদটি প্রথম $200$ ছাড়িয়ে যায়?
\end{enumerate}

\begin{sol}
\begin{enumerate}
  \item $a_{25} = 7 + 24 \times 5 = 127$।
  \item $S_{25} = \dfrac{25}{2}(7 + 127) = \dfrac{25 \times 134}{2} = 1675$।
  \item $7 + (n-1)5 > 200$ চাই, অর্থাৎ $5(n-1) > 193$, অর্থাৎ
        $n - 1 > 38.6$, অর্থাৎ $n > 39.6$। সুতরাং $n = 40$, আর
        $a_{40} = 7 + 39 \times 5 = 202$। আগেরটি $a_{39} = 197$, যা এখনো
        $200$ ছাড়ায়নি।
\end{enumerate}
\end{sol}

\prob{3} একটি সমান্তর অনুক্রমের প্রথম $n$টি পদের যোগফল $S_n = 2n^{2} + 3n$।
$n$-তম পদ বের করো, এবং দেখাও যে অনুক্রমটি সত্যিই সমান্তর।

\begin{sol}
মূল কৌশলটি হলো: $n$-তম পদ হলো প্রথম $n$টির যোগফল বিয়োগ প্রথম $n-1$টির যোগফল,
কারণ মাঝের সব পদ কেটে যায়। অর্থাৎ $a_n = S_n - S_{n-1}$ যখন $n \geq 2$।

\[
  S_{n-1} = 2(n-1)^{2} + 3(n-1) = 2n^{2} - 4n + 2 + 3n - 3 = 2n^{2} - n - 1 .
\]
সুতরাং
\[
  a_n = (2n^{2} + 3n) - (2n^{2} - n - 1) = 4n + 1 .
\]

$n = 1$-এর জন্যও এটি খাটে কি না দেখি: সূত্র দেয় $a_1 = 5$, আর
$S_1 = 2 + 3 = 5$। মিলে গেছে, তাই সূত্রটি সব $n$-এর জন্য বৈধ।

সমান্তর হওয়ার শর্ত হলো পরপর দুই পদের অন্তর ধ্রুবক:
\[
  a_n - a_{n-1} = (4n+1) - \big( 4(n-1) + 1 \big) = 4 ,
\]
যা $n$-এর উপর নির্ভর করে না। সুতরাং অনুক্রমটি সমান্তর, প্রথম পদ $5$ ও সাধারণ
অন্তর $4$।
\end{sol}

\prob{2} $2, 6, 18, 54, \dots$ গুণোত্তর অনুক্রমটির প্রথম $8$টি পদের যোগফল বের
করো।

\begin{sol}
$a = 2$ এবং $r = \tfrac{6}{2} = 3$। সূত্রে
\[
  S_8 = \frac{2\,(3^{8} - 1)}{3 - 1} = 3^{8} - 1 = 6561 - 1 = 6560 .
\]
\end{sol}

\prob{3} $0.4272727\ldots$ দশমিকটিকে ভগ্নাংশে প্রকাশ করো, অসীম গুণোত্তর ধারা
ব্যবহার করে।

\begin{sol}
সংখ্যাটিকে দুই টুকরো করি: যে অংশটি পুনরাবৃত্ত নয়, আর যেটি পুনরাবৃত্ত।
\[
  0.4272727\ldots = 0.4 + \big( 0.027 + 0.00027 + 0.0000027 + \cdots \big) .
\]
বন্ধনীর ভেতরের ধারাটি গুণোত্তর, যার $a = 0.027$ আর
$r = \dfrac{0.00027}{0.027} = 0.01$। যেহেতু $|r| < 1$, সূত্র খাটে:
\[
  S = \frac{0.027}{1 - 0.01} = \frac{0.027}{0.99} = \frac{27}{990}
    = \frac{3}{110} .
\]
সুতরাং
\[
  0.4272727\ldots = \frac{2}{5} + \frac{3}{110}
  = \frac{44}{110} + \frac{3}{110} = \frac{47}{110} .
\]
যাচাই করতে $47$-কে $110$ দিয়ে ভাগ করে দেখো; সত্যিই $0.42727\ldots$ পাওয়া যায়।
\end{sol}

\prob{2} নিচেরগুলো করো।
\begin{enumerate}
  \item $\displaystyle \sum_{i=1}^{4} (2i - 1)^{2}$ খুলে লিখে মান বের করো।
  \item $1 + 4 + 9 + \cdots + 100$ যোগফলটি সিগমা নোটেশনে লেখো।
\end{enumerate}

\begin{sol}
\begin{enumerate}
  \item $i = 1, 2, 3, 4$ বসিয়ে পদগুলো $1^{2}, 3^{2}, 5^{2}, 7^{2}$, তাই
        \[
          1 + 9 + 25 + 49 = 84 .
        \]
  \item পদগুলো $1^{2}$ থেকে $10^{2}$ পর্যন্ত বর্গ, তাই
        $\displaystyle \sum_{i=1}^{10} i^{2}$। (সূত্রে মান
        $\tfrac{10 \times 11 \times 21}{6} = 385$।)
\end{enumerate}
\end{sol}

\prob{3} $\displaystyle \sum_{i=1}^{n} i(i+2)$ নির্ণয় করো, এবং $n = 12$-এর
জন্য মান বের করো।

\begin{sol}
প্রথমে পদটি খুলি, কারণ আমাদের সূত্রগুলো $i$ ও $i^{2}$-এর জন্য:
\[
  \sum_{i=1}^{n} i(i+2) = \sum_{i=1}^{n} (i^{2} + 2i)
  = \sum_{i=1}^{n} i^{2} + 2 \sum_{i=1}^{n} i .
\]
সূত্র বসাই:
\[
  = \frac{n(n+1)(2n+1)}{6} + 2 \cdot \frac{n(n+1)}{2}
  = \frac{n(n+1)(2n+1)}{6} + n(n+1) .
\]
$\dfrac{n(n+1)}{6}$ সাধারণ উৎপাদক নিয়ে
\[
  = \frac{n(n+1)}{6}\Big[ (2n+1) + 6 \Big] = \frac{n(n+1)(2n+7)}{6} .
\]

$n = 12$ বসিয়ে
\[
  \frac{12 \times 13 \times 31}{6} = 2 \times 13 \times 31 = 806 .
\]
\end{sol}

\prob{2} $(x + 2y)^{5}$ বিস্তার করো।

\begin{sol}
প্যাসকেলের পঞ্চম সারি $1, 5, 10, 10, 5, 1$। $a = x$, $b = 2y$ ধরে
\[
  (x + 2y)^{5} = x^{5} + 5x^{4}(2y) + 10x^{3}(2y)^{2} + 10x^{2}(2y)^{3}
               + 5x(2y)^{4} + (2y)^{5} .
\]
$2$-এর ঘাতগুলো হিসাব করে
\[
  = x^{5} + 10x^{4}y + 40x^{3}y^{2} + 80x^{2}y^{3} + 80xy^{4} + 32y^{5} .
\]
লক্ষ করো সহগগুলো আর প্রতিসম নয়, কারণ $b$-তে একটি $2$ লুকিয়ে আছে।
\end{sol}

\prob{3} $\left( x^{2} + \dfrac{2}{x} \right)^{9}$ বিস্তারে $x$-মুক্ত পদটি
বের করো।

\begin{sol}
সাধারণ পদটি লিখি:
\[
  \binom{9}{k} \left( x^{2} \right)^{9-k} \left( \frac{2}{x} \right)^{k}
  = \binom{9}{k}\, 2^{k}\, x^{\,18 - 2k}\, x^{-k}
  = \binom{9}{k}\, 2^{k}\, x^{\,18 - 3k} .
\]
$x$-মুক্ত পদ মানে সূচক শূন্য:
\[
  18 - 3k = 0 \ \Longrightarrow\ k = 6 .
\]
$k = 6$ পরিসরে আছে, তাই এমন পদ সত্যিই একটি আছে। তার মান
\[
  \binom{9}{6}\, 2^{6} = \frac{9!}{6!\,3!} \times 64 = 84 \times 64 = 5376 .
\]
\end{sol}

\prob{3} দ্বিপদী উপপাদ্য ব্যবহার করে নিচেরগুলো দেখাও।
\begin{enumerate}
  \item $\displaystyle \binom{n}{k} = \binom{n}{n-k}$
  \item $\displaystyle \sum_{k=0}^{n} \binom{n}{k} = 2^{n}$
  \item $n \geq 1$ হলে
        $\displaystyle \sum_{k=0}^{n} (-1)^{k} \binom{n}{k} = 0$
\end{enumerate}

\begin{sol}
\begin{enumerate}
  \item সংজ্ঞায় $k$-এর জায়গায় $n-k$ বসাই:
        \[
          \binom{n}{n-k} = \frac{n!}{(n-k)!\,\big( n - (n-k) \big)!}
                         = \frac{n!}{(n-k)!\,k!} = \binom{n}{k} .
        \]
        হরের দুটি উৎপাদক কেবল জায়গা বদল করেছে। এই কারণেই প্যাসকেলের প্রতিটি
        সারি বাঁ থেকে ডানে পড়লে আর ডান থেকে বাঁয়ে পড়লে একই দেখায়।
  \item উপপাদ্যে $a = b = 1$ বসাই। বাঁ পাশ $(1+1)^{n} = 2^{n}$, আর ডান পাশে
        $1^{\,n-k} 1^{\,k} = 1$, তাই
        \[
          2^{n} = \sum_{k=0}^{n} \binom{n}{k} .
        \]
        অর্থাৎ প্যাসকেলের $n$-তম সারির সব সংখ্যা যোগ করলে $2^{n}$ পাওয়া যায়;
        চতুর্থ সারিতে $1+4+6+4+1 = 16 = 2^{4}$।
  \item এবার $a = 1$, $b = -1$ বসাই। বাঁ পাশ $(1 - 1)^{n} = 0^{n}$, যা
        $n \geq 1$ হলে $0$। ডান পাশে $b^{k} = (-1)^{k}$, তাই
        \[
          0 = \sum_{k=0}^{n} (-1)^{k} \binom{n}{k} .
        \]
        অর্থাৎ এক সারির সংখ্যাগুলোতে পর্যায়ক্রমে যোগ-বিয়োগ চিহ্ন বসালে
        যোগফল শূন্য; চতুর্থ সারিতে $1 - 4 + 6 - 4 + 1 = 0$।
        $n = 0$ বাদ রাখতে হলো, কারণ তখন যোগফল $\binom{0}{0} = 1$, শূন্য নয়।
\end{enumerate}
\end{sol}

\end{document}
